\documentclass[../../main.tex]{subfiles}
\graphicspath{{../../image/}{../../figure/}}
\begin{document}

\section{熵耗散诱导的解衰减}\label{section:ch11-熵耗散诱导的解衰减}

真正非线性产生多种耗散机制, 它们单独或共同作用, 影响解的大时间行为。本节将介绍以熵耗散为主要阻尼机制的例子。

\begin{theorem}\label{theorem:ch11-11-5-1}
设 $u$ 是初值为 $u_0$ 的可容许解, 满足
\begin{align}
\int_x^{x+\ell}u_0(y)\,\mathrm{d}y=O(\ell^r),\quad \ell\to\infty,\label{align:ch11-11-5-1}
\end{align}
其中 $r\in[0,1)$ 与 $x\in(-\infty,\infty)$ 一致。则
\begin{align}
u(x\pm,t)=O\left(t^{-\frac{1-r}{2-r}}\right),\quad t\to\infty,\label{align:ch11-11-5-2}
\end{align}
与 $x\in(-\infty,\infty)$ 一致。
\end{theorem}

\begin{proof}
固定 $(x,t)\in(-\infty,\infty)\times(0,\infty)$, 在\eqref{align:ch11-11-4-1}中取 $u^*\equiv0$。注意 $\xi(0)-\xi^*(0)=t[f'(u(x\pm,t))-f'(0)]$。又回忆左端两个积分均非正。于是利用\eqref{align:ch11-11-5-1}, 推出
\begin{align}
\Phi(u(x\pm,t))=O(t^{r-1}),\quad t\to\infty,\label{align:ch11-11-5-3}
\end{align}
与 $x$ 一致, 其中
\begin{align}
\Phi(u)=\frac{f(0)-f(u)+uf'(u)}{|f'(u)-f'(0)|^r}=\frac{\int_0^u vf''(v)\,\mathrm{d}v}{\left|\int_0^u f''(v)\,\mathrm{d}v\right|^r}.\label{align:ch11-11-5-4}
\end{align}
简单估计给出 $\Phi(u)\geqslant K|u|^{2-r}$, $K>0$, 故\eqref{align:ch11-11-5-3}蕴含\eqref{align:ch11-11-5-2}。证毕。
\end{proof}

特别地, 当 $u_0\in L^p$ 时, 由 Hölder 不等式, \eqref{align:ch11-11-5-1} 对 $r=1-\dfrac1p$ 成立。因此\cref{theorem:ch11-11-5-1}有如下推论。

\begin{theorem}\label{theorem:ch11-11-5-2}
设 $u$ 是初值 $u_0\in L^p(-\infty,\infty)$, $1\leqslant p<\infty$ 的可容许解。则
\begin{align}
u(x\pm,t)=O\left(t^{-\frac{p}{p+1}}\right),\quad t\to\infty,\label{align:ch11-11-5-5}
\end{align}
与 $x\in(-\infty,\infty)$ 一致。
\end{theorem}

上述例子中比较函数为解 $u^*\equiv0$。接下来考虑比较函数是由两个常状态 $u_-$ 与 $u_+$, $u_->u_+$, 经激波连接而成的 Riemann 问题解的情形, 即
\begin{align}
u^*(x,t)=
\begin{cases}
u_-, & x<st,\\
u_+, & x>st,
\end{cases}\label{align:ch11-11-5-6}
\end{align}
其中
\begin{align}
s=\frac{f(u_+)-f(u_-)}{u_+-u_-}.\label{align:ch11-11-5-7}
\end{align}

\begin{theorem}\label{theorem:ch11-11-5-3}
设 $u$ 是初值为 $u_0$ 的可容许解, 使得反常积分 $\int_{-\infty}^0[u_0(y)-u_-]\,\mathrm{d}y$ 与 $\int_0^\infty[u_0(y)-u_+]\,\mathrm{d}y$ 存在, 其中 $u_->u_+$。归一化原点 $x=0$ 使得
\begin{align}
\int_{-\infty}^0[u_0(y)-u_-]\,\mathrm{d}y+\int_0^\infty[u_0(y)-u_+]\,\mathrm{d}y=0.\label{align:ch11-11-5-8}
\end{align}
考虑从 $(0,0)$ 出发的任意前向特征线 $\chi(\cdot)$。则当 $t\to\infty$ 时,
\begin{align}
\chi(t)=st+o(1),\label{align:ch11-11-5-9}
\end{align}
其中 $s$ 由\eqref{align:ch11-11-5-7}给出, 且
\begin{align}
u(x\pm,t)=
\begin{cases}
u_-+o(t^{-1/2}), & \text{在 }x<\chi(t)\text{ 上一致},\\
u_++o(t^{-1/2}), & \text{在 }x>\chi(t)\text{ 上一致}.
\end{cases}\label{align:ch11-11-5-10}
\end{align}
\end{theorem}

\begin{proof}
固定任意 $(x,t)\in(-\infty,\infty)\times(0,\infty)$, 对初值为 $u_0$ 的解 $u$ 与由\eqref{align:ch11-11-5-6}给出的比较解 $u^*$ 写出\eqref{align:ch11-11-4-1}。由 $f'(u_-)>s>f'(u_+)$, 当 $t\to\infty$ 时, $\xi^*(0)\to-\infty$ 在 $(-\infty,st)$ 上关于 $x$ 一致, $\xi^*(0)\to\infty$ 在 $(st,\infty)$ 上关于 $x$ 一致。类似地, 当 $t\to\infty$ 时, $\xi(0)\to-\infty$ 在 $(-\infty,\chi(t))$ 上关于 $x$ 一致, $\xi(0)\to\infty$ 在 $(\chi(t),\infty)$ 上关于 $x$ 一致。事实上, 在相反情形, 我们将能找到一列 $(x_n,t_n)$, $t_n\to\infty$, $n\to\infty$, 使得从 $(x_n,t_n)$ 出发的极小后向特征线 $\xi_n(\cdot)$ 的截距 $\xi_n(0)$ 有界。但那样 $\{\xi_n(\cdot)\}$ 的某子列将收敛到从某点 $(\bar x,0)$ 出发的分界线。然而这不可能, 因为 $u_->u_+$, \eqref{align:ch11-11-5-8}与\eqref{align:ch11-11-4-2}对任意 $\bar x\in(-\infty,\infty)$ 和任意 $\bar u\in(-\infty,\infty)$ 不相容。

由上, \eqref{align:ch11-11-5-8}蕴含当 $t\to\infty$ 时, \eqref{align:ch11-11-4-1}右端为 $o(1)$, 在 $(-\infty,\infty)$ 上关于 $x$ 一致。于是左端两个积分各自也为 $o(1)$, 因为它们同号(非正)。

先考虑 $(x,t)\in(-\infty,\infty)\times(0,\infty)$ 且 $x<\min\{\chi(t),st\}$。则 $\xi(\tau)<s\tau$, $0<\tau<t$, 故\eqref{align:ch11-11-4-1}左端第一个积分给出
\begin{align}
t\{f(u(x\pm,t))-f(u_-)-f'(u(x\pm,t))[u(x\pm,t)-u_-]\}=o(1).\label{align:ch11-11-5-11}
\end{align}
由于 $f$ 一致凸, \eqref{align:ch11-11-5-11} 蕴含 $u(x\pm,t)-u_-=o(t^{-1/2})$。

类似论证表明, 对 $(x,t)\in(-\infty,\infty)\times(0,\infty)$ 且 $x>\max\{\chi(t),st\}$, 有 $u(x\pm,t)-u_+=o(t^{-1/2})$。

接下来考虑 $(x,t)\in(-\infty,\infty)\times(0,\infty)$ 且 $st\leqslant x<\chi(t)$。则 $\xi(\cdot)$ 必与直线 $x=s\tau$ 相交, 设在 $\tau=r$, $r\in[0,t]$, 此时\eqref{align:ch11-11-4-1}左端第一个积分给出
\begin{align}
(t-r)\{f(u(x\pm,t))-f(u_+)-f'(u(x\pm,t))[u(x\pm,t)-u_+]\}=o(1),\label{align:ch11-11-5-12}\\
r\{f(u(x\pm,t))-f(u_-)-f'(u(x\pm,t))[u(x\pm,t)-u_-]\}=o(1).\label{align:ch11-11-5-13}
\end{align}

对 $x<\chi(t)$, 如上已证 $\xi(0)\to-\infty$, $t\to\infty$, 这进而蕴含 $r\to\infty$。于是由\eqref{align:ch11-11-5-13}与 $f$ 的凸性, $u(x\pm,t)=u_-+o(1)$。然后\eqref{align:ch11-11-5-12} 蕴含 $t-r=o(1)$, 故 $\chi(t)-st=o(1)$, 而\eqref{align:ch11-11-5-13}给出\eqref{align:ch11-11-5-11}。由\eqref{align:ch11-11-5-11}与 $f$ 的凸性, 如前推出 $u(x\pm,t)-u_-=o(t^{-1/2})$。

类似论证建立: 对 $(x,t)\in(-\infty,\infty)\times(0,\infty)$ 且 $\chi(t)<x\leqslant st$, 有 $u(x\pm,t)-u_+=o(t^{-1/2})$ 且 $\chi(t)-st=o(1)$。证毕。
\end{proof}

\end{document}
